\subsection{Integral formulas for moment ratios and a series lower bound}
\label{subsec:marker-floor}
The lower bound $F_*$ in Lemma~\ref{lem:calibration} contains the
logarithmic function $g$. We bound its contribution using the moments
of $U$ and the successive differences of the ratios $a_s/\beta_s$.
An integral representation relates these differences at all orders.

For real $s\ge0$ retain
\[
 a_s=\E U^s,\qquad \beta_s=\frac{h}{\lambda^s-d},\qquad
 \rho(s)=\frac{a_s}{\beta_s}.
\]
For $s\ge1$ put
\begin{equation}\label{eq:marker-deficits}
 \begin{gathered}
 r_s=\frac{\beta_{s+1}}{\beta_s},\qquad
 \Delta_s=\rho(s)-\rho(s+1),\qquad
 \delta_s=a_s r_s-a_{s+1}=\beta_{s+1}\Delta_s,\\
 r_0=u,\qquad \delta_0=0,\qquad
 k=\frac{c}{\lambda\log\lambda}.
 \end{gathered}
\end{equation}
In particular, $1/\lambda<k<1$.

\begin{lemma}[Integral representation of the moment ratios]\label{lem:canonical-marker}
There is a probability measure $\sigma$ on $(0,1)$, constructed from the
escape law $\pi_\lambda$, such that, for every real $s\ge1$,
\begin{equation}\label{eq:marker-representation}
 \rho(s)=\int_0^1 t^{s-1}\{s-(s-1)t\}\,\sigma(\dd t),
 \qquad
 \Delta_s=s\int_0^1 t^{s-1}(1-t)^2\,\sigma(\dd t).
\end{equation}
The derivative $\rho'(s)$ exists, with the right derivative taken at $s=1$,
and
\begin{equation}\label{eq:strict-marker-derivative}
 -\rho'(s)>\left(1-\frac1{2s}\right)\Delta_s.
\end{equation}
Thus, with $e=\Delta_1=1-\rho_2$, one has $0<e<1$ and
$0<\Delta_i\le i e$ for every integer $i\ge1$.
\end{lemma}
\begin{proof}
Apply the chosen-child inequality~\eqref{eq:chosen-child-comparison} to
increasing indicator functions.  After the same change of variables as
in the proof of Lemma~\ref{lem:canonical-moments}, it says that
\[
 A_m(t)=\E[(1-\lambda U)\mathbf1_{\{\lambda U>t\}}]
       -d\E[(1-U)\mathbf1_{\{U>t\}}]\ge0,
 \qquad 0<t<1.
\]
Layer-cake integration gives
\[
 h\{\rho(s)-\rho(s+1)\}=s\int_0^1t^{s-1}A_m(t)\dd t,
 \qquad s\ge1.
\]
Indeed, the right side equals
$(\lambda^s-d)a_s-(\lambda^{s+1}-d)a_{s+1}$.
The bounded-convergence argument in Lemma~\ref{lem:canonical-moments}
also gives $\rho(s+n)\to0$ for every real $s\ge1$.  Sum the last
display over $s,s+1,\ldots$, using Tonelli and
\[
 \sum_{j\ge0}(s+j)t^{s+j-1}
 =\frac{t^{s-1}\{s-(s-1)t\}}{(1-t)^2}.
\]
At $s=1$ the resulting left side is $\rho(1)=1$. Consequently
\[
 \sigma(\dd t)=\frac{A_m(t)}{h(1-t)^2}\dd t
\]
is a probability measure and gives~\eqref{eq:marker-representation}.
The measure gives full mass to $(0,1)$.

For $s>1$, differentiation under this probability integral is justified
locally by boundedness of $t^{s-1}|\log t|$.  At $s=1$, the derivative
of $a_s/\beta_s$ is finite because $U^s|\log U|$ is bounded when $s$
is bounded away from zero.  Fatou's lemma applied to
\[
 \frac{1-t^{\varepsilon}\{1+\varepsilon(1-t)\}}{\varepsilon}
\]
then shows that $-\log t-(1-t)$ is $\sigma$-integrable.  Hence
$|\log t|$ is integrable and dominated convergence also permits
differentiation from the right at $s=1$. We obtain
\[
 -\rho'(s)=\int_0^1 t^{s-1}
 \bigl[\{1+(s-1)(1-t)\}(-\log t)-(1-t)\bigr]\sigma(\dd t).
\]
Write $z=1-t$.  Since $-\log(1-z)>z+z^2/2$ for $0<z<1$, the
bracket is strictly larger than
$(s-\tfrac12)z^2+(s-1)z^3/2$.
Integration proves~\eqref{eq:strict-marker-derivative}.  Finally,
the formula for $\Delta_i$ and $t^{i-1}\le1$ give $\Delta_i\le i e$.
Strict positivity follows from the support of $\sigma$.
\end{proof}

We now replace the speed $v$ by a real parameter $w$ while keeping the
law of $X$ fixed. Since $wZ-N$ vanishes only at $w=v$, we keep this
term when substituting an upper bound for $v$. In the next two lemmas,
$w$ denotes this parameter. It previously denoted $\E W$
in~\eqref{eq:archive-parameters}.

\begin{lemma}[An affine comparison at fixed escape law]\label{lem:assigned-floor}
Keep the escape law $\pi_\lambda$ and its moments fixed, and put
\[
 A_0=1-(\lambda+1)u+\lambda a_2,\qquad
 a(w)=\frac{u-(w/c)A_0}{a_2},
\]
\[
 P_c(w;y)=(1-y)\{2w-(1+w)y\}.
\]
Define the affine function
\begin{equation}\label{eq:assigned-floor-definition}
 \widetilde F(w)=2a(w)\E[br]
 -\E[b\cJ P_c(w;\cdot)]
 +2\E\!\left[g\cJ\{(1-a(w))y^2+c((1+w)y-w)\}\right].
\end{equation}
If $0<\alpha<1$ in Lemma~\ref{lem:calibration}, then
$\widetilde F(v)=F_*$ and $\widetilde F'(w)<0$ for every real $w$.
In particular, $F_*>\widetilde F(w)$ when $w>v$.
\end{lemma}
\begin{proof}
Set $\Gamma=\E[br]-\E[g\cJ(y^2)]$.  Apply
\eqref{eq:entropy-stationarity} with the strictly positive function $r$,
using $(I-\cL)r=\cJ(y^2)$.  It gives
\[
 \Gamma=\E[\delta_g r(F_\lambda(S))]>0.
\]
All these expectations are finite by~\eqref{eq:entropy-integrability}.
To find the dependence on $w$, symmetrize the bounded
pair kernels.  For $t=x+y-xy$,
\[
 \frac{b(x)P_c(w;y)+b(y)P_c(w;x)}2
 =\frac{(1-t)(wt-xy)}\lambda.
\]
Since $1-t=\lambda-A$, this implies the exact correction
\begin{equation}\label{eq:assigned-C-correction}
 \E\frac{wt-XY}{A^2}-\E[b\cJ P_c(w;\cdot)]
 =\frac{wZ-N}\lambda.
\end{equation}
At $w=v=N/Z$ the correction vanishes, giving
$\widetilde F(v)=F_*$.  Differentiating the affine expression with the
law fixed gives
\[
 \widetilde F'(w)=2c\E[g\cJ(y-1)]
 -\E[b\cJ((1-y)(2-y))]-\frac{2A_0\Gamma}{c a_2}<0.
\]
Here $A_0>0$, $a_2>0$, $\Gamma>0$, $g\ge0$, and
$0<X,Y<1$ almost surely. Only the scalar $w$ varies in this derivative.
\end{proof}

\begin{lemma}[A series lower bound in terms of the moments]\label{lem:mellin-floor}
Fix $w$ satisfying $v\le w<h$ and $0<a(w)<1$.
Set
\begin{align}
 Q_{a(w)}(z)&=-\frac wc+\left(1+(\lambda+1)\frac wc\right)z
                   -\left(a(w)+\lambda\frac wc\right)z^2,
 \nonumber\\
 t_n(w)&=\E[U^nQ_{a(w)}(U)],\label{eq:assigned-t}\\
 J_{s-1}(w)&=-2w a_{s-1}
  +\{c(1+2a(w))+(3\lambda+1)w\}a_s
  \nonumber\\
 &\quad-\{c(\lambda+2a(w))+\lambda(\lambda+1)w\}a_{s+1},
 \label{eq:assigned-J}\\
 \overline\gamma_{s-1}
 &=a_s(r_s-r_{s-1})-\delta_s
   +k\left(1-\frac1{2s}\right)\frac{\delta_s}{r_s}.
 \label{eq:gamma-floor}
\end{align}
Then $t_0(w)=0$, $t_n(w)>0$ for every $n\ge1$, and the series
\begin{equation}\label{eq:complete-mellin-floor}
 F_{\rm floor}(w)=\sum_{s\ge1}s\lambda^{s-1}
 \left\{\frac{(a_s-\lambda a_{s+1})J_{s-1}(w)}{\lambda c}
             +2\overline\gamma_{s-1}t_{s-1}(w)\right\}
\end{equation}
converges absolutely.  Moreover,
\begin{equation}\label{eq:strict-Mellin-gap}
 \widetilde F(w)-F_{\rm floor}(w)
 =2\sum_{s\ge2}s\lambda^{s-1}
 \left\{k\beta_s[-\rho'(s)]
 -k\left(1-\frac1{2s}\right)\frac{\delta_s}{r_s}\right\}
 t_{s-1}(w)>0.
\end{equation}
In particular, $D_\lambda\ge F_*>F_{\rm floor}(v)$.
\end{lemma}
\begin{proof}
The choice of $a(w)$ gives $t_0(w)=0$. The quadratic $Q_{a(w)}$ is negative
at zero, strictly concave, and has value $(\lambda-a(w))/\lambda^2>0$
at $1/\lambda$.  It therefore has exactly one zero $z_0$ in
$(0,1/\lambda)$, with the sign of $z-z_0$ on that interval.  For
$n\ge1$,
\[
 t_n(w)=\E[(U^n-z_0^n)Q_{a(w)}(U)]>0.
\]
The integrand is nonnegative and vanishes only at $U=z_0$.
The tree starts with a chain of single-child vertices of any given
length with positive probability. Along such a chain of length $j$,
the transform $U$ is multiplied by $\lambda^{-j}$. Thus the
essential infimum of $U$ is zero, although $U>0$ almost surely. In particular, its law cannot be concentrated at $z_0>0$.

Write $V=Y/(c+Y)$ for an independent copy of $U$.  Direct substitution
gives, with $a=a(w)$,
\begin{align*}
 \frac{b(X)Y}{A^2}
 &=\frac{UV(1-\lambda U)(1-V)}{\lambda(1-\lambda UV)^2},\\
 \frac{b(X)P_c(w;Y)}{A^2}
 &=\frac{U(1-\lambda U)(1-\lambda V)
            \{2w-(c+(\lambda+1)w)V\}}
          {\lambda c(1-\lambda UV)^2},\\
 \frac{(1-a)Y^2+c((1+w)Y-w)}{A^2}
 &=\frac{(1-U)^2Q_a(V)}{(1-\lambda UV)^2}.
\end{align*}
Since $\lambda UV\le1/\lambda<1$, expand the common denominator in
its uniformly convergent series.  With
\[
 \gamma_{s-1}=\E[g(X)(1-U)^2U^{s-1}]
 =k\E[U^s\log(1/U)]-a_s+a_{s+1},
\]
the first two kernels produce
$(a_s-\lambda a_{s+1})J_{s-1}(w)/(\lambda c)$, and the third produces
$2\gamma_{s-1}t_{s-1}(w)$.  Their sum is exactly
$\widetilde F(w)$.  In particular, the first factor is
$a_s-\lambda a_{s+1}$ for every index $s\ge1$, including the first.

Differentiate $a_s=\beta_s\rho(s)$ with respect to the moment order $s$.
Since
$-k\beta_s'/\beta_s=1-r_{s-1}$, one obtains
\[
 \gamma_{s-1}=a_s(r_s-r_{s-1})-\delta_s
                    +k\beta_s[-\rho'(s)].
\]
Use Lemma~\ref{lem:canonical-marker},
$\beta_s\Delta_s=\delta_s/r_s$, and $t_{s-1}(w)>0$ for $s\ge2$.
The term $s=1$ has $t_0(w)=0$.  This proves the equality and strict
inequality in~\eqref{eq:strict-Mellin-gap}.

For completeness, $a_s\le\lambda^{-s}$ and
$\delta_s\le\lambda^{-s-1}$.  The ratios $r_s$ are bounded away from
zero at fixed $\lambda,h$, and $r_s-r_{s-1}$ is bounded.  The integral
defining $\gamma_{s-1}$ is bounded by
$\|g\|_\infty\lambda^{-s+1}$; its identity in the preceding display
therefore bounds $k\beta_s[-\rho'(s)]$ by a constant times
$\lambda^{-s}$.  Also $t_{s-1}(w)$ and $J_{s-1}(w)$ have bounds of
the same geometric order.  Every expansion and regrouping above is
dominated by a constant times $\sum_{s\ge1}s\lambda^{-s}$.
No unweighted logarithmic moment or inverse escape moment is used.
\end{proof}
